(1) 三维自由电子满足 \(\epsilon=p^2/(2m)\)，因此态密度 \(g(\epsilon)=A\sqrt{\epsilon}\)，其中 \(A\) 已含自旋两重简并。利用 \(N=\int_0^{E_F}A\sqrt\epsilon\,d\epsilon=(2/3)AE_F^{3/2}\)，得 \(A=3N/(2E_F^{3/2})\)。
(2) Sommerfeld展开为 \(\int f(\epsilon)n_F(\epsilon)d\epsilon=\int_0^\mu f(\epsilon)d\epsilon+\pi^2(k_BT)^2f'(\mu)/6+O(T^4)\)。将 \(f=g\) 代入固定 \(N\) 条件，解出 \(\mu(T)=E_F[1-(\pi^2/12)(k_BT/E_F)^2+O(T^4)]\)。
(3) 对 \(U=\int_0^\infty\epsilon g(\epsilon)n_F(\epsilon)d\epsilon\) 展开并代入修正化学势，得到 \(U=(3/5)NE_F+(\pi^2/4)N(k_BT)^2/E_F+O(T^4)\)。因此 \(C_V=(\pi^2/2)Nk_B(k_BT/E_F)+O(T^3)\)，在低温按 \(T\) 线性增长。
(4) 经典单原子气体有 \(C_V=3Nk_B/2\)，但简并电子气低温满足 \(C_V/(Nk_B)\simeq(\pi^2/2)T/T_F\)，其中 \(T_F=E_F/k_B\)。因Fermi海深处能级已占满，只有能量距Fermi面约 \(k_BT\) 的电子容易参与激发。
